\section*{Resumen}

Se desarrolló el diseño sísmico de los muros Pier~C y Pier~J y de las vigas sísmicas B7, B8 y B39
de un edificio de hormigón armado de 21 pisos ubicado en zona sísmica~3, suelo tipo~A. Los muros,
de espesor $e_w = 25\;\text{cm}$ y altura total $H_w = 50{,}82\;\text{m}$, se diseñaron con
armadura de corte horizontal con sobrerresistencia $\Omega_v = 2{,}0$ y armadura de
flexo-compresión mediante el método de Cárdenas-Magura, adoptando $\rho_v = 0{,}0035$ en todos
los niveles. El desplazamiento último de demanda resulta $\delta_u = 19{,}57\;\text{cm}$, con
curvaturas últimas de $\phi_u = 1{,}427 \times 10^{-3}\;\text{m}^{-1}$ (Pier~C) y
$\phi_u = 1{,}337 \times 10^{-3}\;\text{m}^{-1}$ (Pier~J). El único piso que requiere
confinamiento especial de borde es el piso~1 del Pier~C, con longitud de confinamiento
$l_\text{conf} = 31{,}4\;\text{cm}$ y separación máxima $s_\text{máx} = 7{,}2\;\text{cm}$.
Para las vigas sísmicas, la armadura longitudinal varía entre $2\phi22$ en pisos altos y $3\phi25$
en pisos bajos, con estribos $\phi8$ a $\phi12$ según la demanda de corte. Todos los elementos
verifican los requisitos de resistencia conforme a ACI~318-19, NCh433 y DS~60.